\section{The discrete Fourier transform}

\begin{correction}{exo-bit-reversal}
\begin{enumerate}
\item The first vectors are
$u^{(1)}=(0,1)$, $u^{(2)}=(0,2,1,3)$, and
$u^{(3)}=(0,4,2,6,1,5,3,7)$.
\item When the most significant input bit is zero, reversal appends a zero
to the reversed lower bits, giving $2u^{(n-1)}$. When it is one, reversal
appends a one, giving $2u^{(n-1)}+1$. This is exactly the defining recursion.
\item In this order, the leaves of the DIT recursion are adjacent in memory.
The iterative algorithm combines contiguous blocks from size two upward;
Section~\ref{sect2-impl-concrete} explains why this evaluates the DFT of the
original signal, rather than the DFT of a different signal.
\item A first-half output index has leading bit zero, so its reversed index
is even. Reversing its remaining bits gives $\wt f_g=\widetilde{f^0}$.
The same argument with a leading one gives $\wt f_d=\widetilde{f^1}$.
\item Program~\ref{listing-revbitsrec} implements this relation.
At every level, concatenation copies $N$ entries in total, so this recursive
array construction costs $O(N\log N)$ element moves.
In contrast, the direct index recursion creates arrays of lengths
$1,2,4,\ldots,N$, for $O(N)$ total work, followed by a single $O(N)$ gather.
The two implementations do not have the same copying complexity.
\begin{listing}[H]
\PythonCode{listing-revbitsrec}
\caption{Recursive permutation of a vector into bit-reversed order.}
\label{listing-revbitsrec}
\end{listing}
\end{enumerate}
\end{correction}

\begin{correction}{exo-algo-good-thomas}
\begin{enumerate}
\item Choose integers $u,v$ with $up+vq=1$. Then
$\psi(k_1,k_2)=(k_1vq+k_2up)\bmod N$ has the required two residues.
The Chinese remainder theorem makes it the inverse of $\varphi$.
\item Since $n=\psi(k_1,k_2)$ satisfies $n\equiv k_1\pmod p$ and
$n\equiv k_2\pmod q$,
\[
 n(s_1q+s_2p)\equiv k_1s_1q+k_2s_2p\pmod N.
\]
This follows either by reducing modulo both coprime factors or by expanding
$\psi$. Substitute in the two-dimensional Fourier sum and reindex by $n$:
\[
 \wh F[s_1,s_2]
 =\sum_{n=0}^{N-1}f[n]\omega_N^{-n(s_1q+s_2p)}
 =\wh f[(s_1q+s_2p)\bmod N].
\]
\item For a given output index $m$, the inverse mapping is
$s_1=mq^{-1}\bmod p$ and $s_2=mp^{-1}\bmod q$.
These inverses exist by coprimality, proving bijectivity.
The integer $s_1q+s_2p$ can exceed $N-1$; the reduction modulo $N$ is essential.
\item The character pairing factors with no additional phase, so there is
no analogue of $\omega_N^{-bd}$ in \eqref{eqn-cooley-tukey-twiddle}.
Program~\ref{listing-fftgt} scatters the input by its two residues, applies
a separable two-dimensional FFT, and scatters the output by the proved map.
Further factoring the row and column transforms combines this method with
Cooley--Tukey or other algorithms. The permutations still cost memory traffic;
removing multiplications does not by itself establish a running-time advantage.
\begin{listing}[H]
\PythonCode{listing-fftgt}
\caption{A Good--Thomas factorization with explicit input and output permutations.}
\label{listing-fftgt}
\end{listing}
\end{enumerate}
\end{correction}

\begin{correction}{exo-algo-split-radix}
\begin{enumerate}
\item Write the input index as $a+bN/4$. Expanding the exponent for
frequency $4k+2j+1$ gives
\[
 \omega_N^{-(a+bN/4)(4k+2j+1)}
 =\omega_{N/4}^{-ak}\omega_N^{-a(2j+1)}\omega_4^{-b(2j+1)},
\]
since the omitted factor $e^{-2\pi\imath bk}$ is one.
With $u_a=f[a]-f[a+N/2]$ and $v_a=f[a+N/4]-f[a+3N/4]$, the two inner
sums are $u_a-\imath v_a$ and $u_a+\imath v_a$.
Their general twiddle factors are $\omega_N^{-a}$ and $\omega_N^{-3a}$.
The even-output branch is the length-$N/2$ DFT of $f_g+f_d$, without twiddles.
\item Let $E$ be the length-$N/2$ DFT of the even input entries, and let
$O_1,O_3$ be the length-$N/4$ DFTs of the entries at indices $4j+1$ and
$4j+3$. For $0\leq k<N/4$, put
$u=\omega_N^{-k}O_1[k]$, $v=\omega_N^{-3k}O_3[k]$.
Then
\begin{align*}
 \wh f[k]&=E[k]+u+v,&
 \wh f[k+N/2]&=E[k]-u-v,\\
 \wh f[k+N/4]&=E[k+N/4]-\imath(u-v),&
 \wh f[k+3N/4]&=E[k+N/4]+\imath(u-v).
\end{align*}
These follow by substituting the four frequency quarters into the three
input-subsequence sums. The function \texttt{split\_radix} in the accompanying
Python module follows these formulas; sizes one and two are the base cases.
\item Let $M(N)$ count the displayed general complex multiplications.
The split-radix combine stage uses $N/2$ such products, giving
\[
 M(N)=M(N/2)+2M(N/4)+N/2,\qquad M(1)=M(2)=0.
\]
For $N=2^p$, solving the recurrence yields
$M(N)=Np/3-2N/9+(2/9)(-1)^p$.
With the same length-two base case, radix two gives $Np/2-N/2$.
Thus the leading coefficient drops from $1/2$ to $1/3$ in this model.
Each split-radix quartet uses six complex additions after sharing $u+v$
and $u-v$, so
\[
 A(N)=A(N/2)+2A(N/4)+3N/2,\qquad A(1)=0,\quad A(2)=2.
\]
It follows that $A(N)=N\log_2N$, the same count as radix two.
Further savings are possible by special-casing trivial twiddles; these were
excluded from this count. An iterative version must schedule subproblems of
two different sizes and track where the three subsequences and four output
quarters are stored. Its permutations must follow that factorization, rather
than assuming equal-length blocks at every stage.
\end{enumerate}
\end{correction}

\begin{correction}{exo-calcul-convolution-fft}
\begin{enumerate}
\item As zero-extended sequences,
$f=\sum_jf_j[\cdot-jM]$. Bilinearity and translation covariance of linear
convolution give the asserted sum of shifted block convolutions.
\item Each block output has length at most $2M-1$.
Choose one FFT length $L\geq2M-1$, for example the next power of two,
and compute the padded transform of $g$ once.
For every block, take one forward FFT, multiply by this stored transform,
and take an inverse FFT. Add the first $2M-1$ entries at offset $jM$.
With $L=O(M)$, each block costs $O(M\log M)$, and there are $N/M$ blocks.
Including the overlap additions gives $O(N\log M)$ for $M\geq2$.
For a one-sample kernel the operation is simply scalar multiplication.
\item If the last block has $r<M$ entries, pad it to the same FFT length
but add only its $r+M-1$ output entries. No output beyond index $N+M-2$
should be retained. Program~\ref{listing-convol} handles this case directly.
\begin{listing}[H]
\PythonCode{listing-convol}
\caption{Overlap-add convolution, including an incomplete final block.}
\label{listing-convol}
\end{listing}
\end{enumerate}
\end{correction}

\begin{correction}{exo-circulating-matrix}
\begin{enumerate}
\item Since $(Rx)[j]=x[j-1]$ modulo $N$,
\[
 (\Omega_NRx)[k]=\sum_jx[j-1]\omega_N^{-kj}
 =\omega_N^{-k}\sum_jx[j]\omega_N^{-kj}=(D\Omega_Nx)[k].
\]
Multiply on the right by $\Omega_N^{-1}$.
\item The first column of $R^j$ is $e_j$, and every column is a cyclic
shift of it; hence $C=\sum_{j=0}^{N-1}c_jR^j$.
Conjugating by $\Omega_N$ gives $\sum_jc_jD^j$.
Its $k$th diagonal entry is $\sum_jc_j\omega_N^{-kj}=\wh c[k]$.
Applying this diagonalization to $x$ gives the multiplication formula.
\item Directly,
$(Cx)[i]=\sum_jc[i-j]x[j]=(c*x)[i]$.
The DFT of this vector is $\wh c\cdot\wh x$, so inversion yields the same result.
\end{enumerate}
\end{correction}

\begin{correction}{exo-interpolation-trigonometric}
A negative signed frequency $k\in\{-N_0,\ldots,-1\}$ is stored at index
$N+k$ in the original spectrum and at index $P+k$ in the padded spectrum.
Both indices represent the same exponential after restriction to the
appropriate grid. Thus inversion gives
\[
 f_0[j]=\frac{\eta}{P}\sum_{k=-N_0}^{N_0}\wh f[k\bmod N]
                         e^{2\pi\imath kj/P}=p(j/P).
\]
At $j=\eta m$, $j/P=m/N$, so DFT inversion gives $p(m/N)=f[m]$.
The factor $\eta=P/N$ compensates for the inverse transform's new divisor.
The result is the unique interpolant with signed frequencies $-N_0,\ldots,N_0$.
Program~\ref{listing-interp-trigo} preserves complex input; it removes only
roundoff-level imaginary parts when the input is real.
\begin{listing}[H]
\PythonCode{listing-interp-trigo}
\caption{Trigonometric interpolation by correctly scaled spectral padding.}
\label{listing-interp-trigo}
\end{listing}
\end{correction}

\begin{correction}{exo-interpolation-chebyshev}
\begin{enumerate}
\item The identity
$\cos((k+1)\theta)+\cos((k-1)\theta)=2\cos\theta\cos(k\theta)$
gives the recurrence. Induction proves the degree assertion, with leading
coefficient $2^{k-1}$ for $k\geq1$.
Substitution gives $T_N(x_j)=\cos((j+1/2)\pi)=0$.
These are $N$ distinct points, so they exhaust the roots.
\item The polynomials $T_0,\ldots,T_{N-1}$ have distinct degrees and nonzero
leading coefficients. They form a basis of the polynomials of degree at most
$N-1$, giving the unique expansion of the unique interpolant.
\item Put $\theta_j=(j+1/2)\pi/N$. A finite geometric sum, or pairing
conjugate exponentials, shows
\[
 \sum_{j=0}^{N-1}\cos(k\theta_j)\cos(l\theta_j)
 =\begin{cases}N,&k=l=0,\\N/2,&k=l>0,\\0,&k\neq l,
 \end{cases}
 \qquad 0\leq k,l<N.
\]
Thus the cosine rows, scaled by $N^{-1/2}$ for row zero and $(2/N)^{1/2}$
for the other rows, form an orthonormal matrix. This proves
$\Cc_2^{-1}=(2/N)\Cc_3$.
For $\Cc_2$, place $v[k]$ at index $2k+1$ of a length-$4N$ zero array.
The real part of its DFT at index $j$ is exactly $\Cc_2(v)[j]$.
For $\Cc_3$, pad $v$ at the end to length $4N$, take the real DFT values
at indices $2j+1$, and subtract $v[0]/2$.
Programs~\ref{listing-dct2} and~\ref{listing-dct3} implement these two embeddings.
\begin{listing}[H]
\PythonCode{listing-dct2}
\caption{The cosine transform $\mathcal C_2$ via a length-$4N$ FFT.}
\label{listing-dct2}
\end{listing}
\begin{listing}[H]
\PythonCode{listing-dct3}
\caption{The cosine transform $\mathcal C_3$ via a length-$4N$ FFT.}
\label{listing-dct3}
\end{listing}
\item The interpolation conditions are
\begin{equation}
\label{eqn-calcul-coef-cheby}
 v[j]=\sum_{k=0}^{N-1}\alpha_k\cos(k\theta_j)
      =\Cc_3(\alpha)[j]+\frac{\alpha_0}{2}.
\end{equation}
Orthogonality gives
$\alpha_0=\Cc_2(v)[0]/N$ and
$\alpha_k=2\Cc_2(v)[k]/N$ for $k>0$.
The coefficients cost $O(N\log N)$ to compute. Evaluating the degree-$N-1$
series at each new point costs $O(N)$ using a recurrence, as in NumPy's
\texttt{chebval}. Program~\ref{listing-interpolation-chebyshev} includes the
Runge example from the exercise.
\begin{listing}[H]
\PythonCode{listing-interpolation-chebyshev}
\caption{Chebyshev interpolation and the Runge example.}
\label{listing-interpolation-chebyshev}
\end{listing}
\end{enumerate}
The Chebyshev nodes cluster near the endpoints, reducing the large endpoint
oscillations seen with equally spaced nodes. Both constructions are polynomial
interpolation; the difference is the node distribution, not whether Lagrange
interpolation is being used.
\end{correction}

\begin{correction}{exo-derivation-fractionionnaire}
Integration by parts gives $\wh{f'}(\xi)=\imath\xi\wh f(\xi)$, since the
boundary term vanishes for a Schwartz function. Iterating gives
$\wh{f^{(m)}}=(\imath\xi)^m\wh f$.
The proposed multiplier extends these integer powers by fixing arguments
$\pi/2$ and $-\pi/2$ on the positive and negative frequency axes.
Its values at opposite frequencies are complex conjugates. Since real input
also has conjugate-symmetric Fourier coefficients, the output is real.
For noninteger $\alpha$, this is a particular fractional differentiation
operator; choosing a different multiplier would define a different operator.

For a period $L$, the discrete frequencies are $2\pi k/L$ with
$k=-(N-1)/2,\ldots,(N-1)/2$ for odd $N$.
Program~\ref{listing-derfrac} multiplies the DFT by the specified powers and
inverts it. At integer orders this differentiates the trigonometric interpolant
exactly in exact arithmetic. The odd sample count ensures that each nonzero
frequency has a distinct negative partner, avoiding an ambiguous Nyquist slot.
The case $\alpha=0$ is explicitly the identity, including on the mean.
\begin{listing}[H]
\PythonCode{listing-derfrac}
\caption{A periodic fractional derivative with explicit order and period.}
\label{listing-derfrac}
\end{listing}
For a rapidly decaying function on $\RR$, the period must be large enough to
make interactions with periodic copies small, and the spacing fine enough to
resolve its Fourier content. A Gaussian is not strictly bandlimited; its decay
allows an approximation, not an exact appeal to Shannon's theorem.
Positive derivative order also amplifies high-frequency noise and roundoff.
Negative orders require a separate treatment of the zero frequency and are
outside the definition in this exercise.
\end{correction}

\begin{correction}{exo-transforme-partial-fourier}
\begin{enumerate}
\item Unitarity implies normality. Since $U^2f=f^\sharp$ and reversal is an
involution, $U^4=I$. Its spectrum is contained in $\{1,-1,\imath,-\imath\}$.
The spectral theorem for normal matrices gives an orthonormal eigenbasis.
The DFT matrix is symmetric but generally not Hermitian.
\item The orthogonal projectors satisfy
$P_\lambda P_\mu=\delta_{\lambda\mu}P_\lambda$ and $\sum_\lambda P_\lambda=I$.
Therefore
\[
 U^\alpha U^\beta=\sum_\lambda
 e^{\imath(\alpha+\beta)\theta_\lambda}P_\lambda=U^{\alpha+\beta}.
\]
The multipliers have modulus one, proving unitarity, and
$e^{\imath\theta_\lambda}=\lambda$ gives $U^1=U$.
A fixed real phase is a consistent logarithm on each eigenspace;
independently chosen roots need not have this additive dependence on the exponent.
\item The matrix $J$ of cyclic reversal has entries
$J_{ij}=1$ when $i+j\equiv0\pmod N$, and zero otherwise.
Thus $U^0=I$ and $U^2=J$: index zero stays fixed, and the other nonzero
entries lie along the reversed diagonal. The intermediate matrices combine
these two sparse patterns with $U$ and $U^3$ through the projector formulas.
For the unnormalized transform,
$\Omega_N^\alpha=N^{\alpha/2}U^\alpha$ and $\Omega_N^2=NJ$.
\item When the same phase is assigned throughout an eigenspace, its
contribution is the scalar $e^{\imath\alpha\theta_\lambda}$ times its
orthogonal projector, which is basis-independent.
Replacing $\theta_\lambda$ by $\theta_\lambda+2\pi m_\lambda$ for fixed
integers $m_\lambda$ produces different families at noninteger exponents,
while preserving the integer powers and the composition law.
The program exposes these optional integer branch shifts.
One could also choose different logarithm branches along separately chosen
orthogonal directions within a repeated eigenspace; that is an additional
choice, not an automatic consequence of diagonalization.
\end{enumerate}
Program~\ref{listing-tfdinterm} uses
$P_\lambda=\frac14\sum_{j=0}^3\lambda^{-j}U^j$ and cyclic reversal.
It avoids relying on a general eigenvalue routine to return orthogonal vectors
inside repeated eigenspaces. Its projector computation needs only one FFT,
since $U^2$ is reversal and $U^3$ is the reversal of $U$.
\begin{listing}[H]
\PythonCode{listing-tfdinterm}
\caption{Unitary intermediate DFTs from explicit spectral projectors.}
\label{listing-tfdinterm}
\end{listing}
\end{correction}

\begin{correction}{exo-diagonalization-tfd}
\begin{enumerate}
\item This follows from the real symmetric spectral theorem.
\item The circulant matrix $\Gamma$ has first column
$(-2,1,0,\ldots,0,1)$. Its DFT is
$2\cos(2\pi k/N)-2$, the diagonal of $\Delta$.
The convolution theorem gives $\Omega_N\Gamma=\Delta\Omega_N$.
Taking transposes, using symmetry of all three matrices, gives
$\Gamma\Omega_N=\Omega_N\Delta$.
Adding shows $\Omega_NS=S\Omega_N$.
\item If $A$ and $B$ commute, each eigenspace of $A$ is invariant under $B$.
The restriction of a diagonalizable operator to an invariant subspace is
diagonalizable: its minimal polynomial divides that of the operator, which
has distinct linear factors. Diagonalize $B$ on each eigenspace of $A$ and
combine the resulting bases. This proves simultaneous diagonalizability.
It does not assert that an arbitrary eigenbasis chosen for $A$ already
consists of eigenvectors of $B$.
\item Cyclic reversal commutes with $S$: the two shifts are interchanged
and the diagonal cosine is even. Thus the two parity subspaces are invariant.
In the bases specified in the exercise, the shifts connect only neighboring
indices $k$ and $k+1$. Their interior off-diagonal entries are one; coupling
to $e_0$ or, for even $N$, to $e_{N/2}$ gives $\sqrt2$ in the symmetric
block. When $N$ is odd, coupling across the last pair contributes $+1$ or
$-1$ to the final diagonal entry of the symmetric or antisymmetric block.
Hence both blocks are tridiagonal, with every existing subdiagonal entry
nonzero. There is a zero connection between the two blocks in the full
transformed matrix; it is not an irreducible tridiagonal matrix.
\item Let $T$ be a real symmetric tridiagonal matrix with subdiagonal
entries $b_i\neq0$. In $(T-\lambda I)v=0$, the first row determines $v_2$
from $v_1$, and each successive row determines $v_{i+1}$ from the preceding
two coordinates. If $v_1=0$, every coordinate is zero.
Thus each eigenspace has dimension at most one. Symmetry gives a full
eigenbasis, so the eigenvalues are distinct. The condition concerns the
\emph{subdiagonal}, not the diagonal entries.
It applies separately to each parity block, whose spectra may overlap.
For $N=4$, the independent vectors $(0,1,0,-1)$ and $(1,-1,-1,-1)$ both
satisfy $Sv=-4v$, explicitly disproving simplicity of the full spectrum.
\item The Fourier matrix also preserves parity and commutes with $S$.
It therefore preserves every one-dimensional eigenspace of each restricted
block. The union of the two real orthonormal block eigenbases is consequently
a common eigenbasis, even at eigenvalues shared between the blocks.
\end{enumerate}
Program~\ref{listing-vectproprestfd} uses this construction. It orders vectors
by decreasing eigenvalue of $S$, resolving numerically equal values by putting
the symmetric block first. It fixes each sign by making a largest-magnitude
coordinate positive, using the first such coordinate in a tie.
The rounding used solely for sorting prevents tiny roundoff differences from
reordering an exact tie. The eigenvectors themselves are not rounded.
\begin{listing}[H]
\PythonCode{listing-vectproprestfd}
\caption{A common eigenbasis obtained by diagonalizing the two parity blocks.}
\label{listing-vectproprestfd}
\end{listing}
\end{correction}

\begin{correction}{exo-orthogonalization-fourier}
\begin{enumerate}
\item Substitute $r=x-k$ in the translation sum to obtain
\[
 \Ff(k\top f)[t]=\sum_xf[x-k]\omega_n^{-tx}
 =\omega_n^{-tk}\wh f[t]=(k\bot\wh f)[t].
\]
Multiplying $f[x]$ by $\omega_n^{-kx}$ instead gives
$\wh f[t+k]=((-k)\top\wh f)[t]$. The sign in the second identity is essential.
\item The unitary transform maps each member of the translation orbit to
the corresponding modulation of $\Ff f/\sqrt n$, preserving all inner
products. In the reverse direction, the minus sign merely permutes the orbit
indices. Using the unnormalized DFT without its scaling would multiply
squared norms by $n$.
\item Write $\wt f[x]=\ol{f[-x]}$. Then
\[
 \dotp f{k\top f}=\frac1n\sum_xf[x]\ol{f[x-k]}
                =\frac1n(f*\wt f)[k].
\]
Translation is unitary, so the whole orbit is orthonormal precisely when
$f*\wt f=n\delta_0$. Since $\wh{\wt f}=\ol{\wh f}$, the convolution
theorem makes this equivalent to $|\wh f[k]|^2=n$ for all $k$.
\item The prescribed normalization gives exactly this constant magnitude.
For modulations, the inner products are the Fourier coefficients of $|f|^2$
with normalized sums, so the orbit is orthonormal precisely when $|f[x]|=1$
for every $x$. Thus $f[x]/|f[x]|$ is the analogous construction when no
sample vanishes. Very small denominators make either phase normalization
sensitive to perturbations, even when exact zeros are absent.
\item The same correlation calculation gives
\[
 \Gg(\varphi)=\frac1n\Ff^{-1}(\wh\varphi\cdot\ol{\wh g}).
\]
It costs $O(n\log n)$ and its output consists of the coefficients in the
orthonormal basis of translates. Therefore
$\varphi=\sum_k\Gg(\varphi)[k](k\top g)=\Gg(\varphi)*g$.
\end{enumerate}
Program~\ref{listing-fftorthog} implements the translation normalization with
the required factor $\sqrt n$ and preserves complex-valued input.
\begin{listing}[H]
\PythonCode{listing-fftorthog}
\caption{Orthonormalizing cyclic translates for the normalized inner product.}
\label{listing-fftorthog}
\end{listing}
\end{correction}
